\section{Approach and Implementation}
\label{sec:approach}

\begin{figure*}[t]
\centering
\fbox{\parbox{0.95\textwidth}{\centering
\textbf{[Pipeline Architecture Overview]}\\[0.5em]
\texttt{IAM JSON Policies} $\longrightarrow$ \texttt{Normalization \& AST Extraction} $\longrightarrow$ \texttt{Heterogeneous Graph Construction} \\
$\longrightarrow$ \texttt{Leak-Free Train/Val Split} $\longrightarrow$ \texttt{Node2Vec Embeddings} $\longrightarrow$ \texttt{Anomaly Scoring (LOF)} $\longrightarrow$ \texttt{Validation Calibration} $\longrightarrow$ \texttt{Locked Test Evaluation}
}}
\caption{End-to-end architecture of the graph-embedding anomaly detection pipeline.}
\label{fig:pipeline}
\end{figure*}

Our architecture translates declarative IAM access policies into low-dimensional structural vectors to detect privilege anomalies without expert-crafted rules. Figure~\ref{fig:pipeline} illustrates the end-to-end pipeline.

\subsection{Heterogeneous IAM Graph Construction}
We formalize cloud IAM authorization topology as a directed heterogeneous graph $G = (V, E)$, consisting of four node types and three directed relation types:
\begin{itemize}
    \item \textbf{Entity Nodes ($V_E$)}: Represent identities, partitioned into Users ($u$), Groups ($g$), and Roles ($r$).
    \item \textbf{Policy Nodes ($V_P$)}: Represent individual IAM managed or inline policies ($p$).
    \item \textbf{Action Nodes ($V_A$)}: Represent standardized, lowercased AWS API actions (e.g., \texttt{"s3:getobject"}, \texttt{"ec2:describesecuritygroups"}). Glob patterns are expanded where finite or retained as explicit wildcard action nodes.
    \item \textbf{Resource Nodes ($V_R$)}: Represent resource patterns, parameterized to normalize organization-specific account IDs and region prefixes.
\end{itemize}

Edges represent authorization dependencies:
\begin{itemize}
    \item $(e, \text{ATTACHED\_TO}, p)$ connects an Entity $e \in V_E$ to an attached Policy $p \in V_P$.
    \item $(p, \text{PERMITS}, a)$ connects Policy $p \in V_P$ to an allowed Action $a \in V_A$.
    \item $(p, \text{TARGETS}, r)$ connects Policy $p \in V_P$ to a target Resource $r \in V_R$.
\end{itemize}
The graph is loaded into Neo4j via Cypher transactions with a fallback to NetworkX for standalone reproduction.

\subsection{Leakage-Safe Node2Vec Embedding Projection}
A major methodological hazard in evaluating graph-embedding anomaly detection is transductive data leakage: if test policies (including injected misconfigurations) are present in the graph when random walks are generated, the resulting embeddings encode relational artifacts of the test distribution. 

To ensure complete leakage safety:
\begin{enumerate}
    \item The graph $G_{\text{train}}$ is constructed exclusively from normal policies in the training split.
    \item Node2Vec biased random walks ($p=1.0, q=1.0$, walk length 80, 10 walks per node, vector dimension $d=128$) are generated strictly over $G_{\text{train}}$, training a Skip-gram model with negative sampling.
    \item Test policies are projected into the learned vector space without retraining. A test policy node $p_{\text{test}}$ inherits its vector by aggregating the embeddings of its incident actions, resources, and entity contexts learned during training, preventing any test-time structural contamination.
\end{enumerate}

\subsection{Anomaly Detection and Calibrated Decision Thresholding}
Once policy embeddings $\mathbf{z}_p \in \mathbb{R}^{128}$ are produced, an unsupervised anomaly detector computes an anomaly score $s(p) \in \mathbb{R}^+$. While our primary detector is Local Outlier Factor ($k=20$), our pipeline supports interchangeable estimators:
\begin{itemize}
    \item \textbf{LOF}: Density-based scoring comparing local reachability against $k$-nearest training neighbors.
    \item \textbf{One-Class SVM}: Non-linear boundary estimation with an RBF kernel ($\gamma = \text{"scale"}$).
    \item \textbf{Isolation Forest}: Ensemble of 100 random partition trees scoring isolation depth.
    \item \textbf{Robust Covariance}: Minimum Covariance Determinant (MCD) estimating Mahalanobis distances from the training centroid.
\end{itemize}

To map continuous anomaly scores $s(p)$ to binary classification decisions ($\hat{y} \in \{0, 1\}$) without leaking test labels:
\begin{enumerate}
    \item Detectors are fitted exclusively on normal training embeddings.
    \item The decision threshold $\tau$ is calibrated using an internal validation split $\mathcal{D}_{\text{val}}$. We select $\tau$ at the $(1 - \alpha)$-th percentile of normal validation scores (e.g., $\alpha = 0.05$ or tuned on a held-out validation tuning grid).
    \item Policies in the locked test set $\mathcal{D}_{\text{test}}$ are predicted as anomalous if $s(p) \ge \tau$.
\end{enumerate}

